\section{Batasan dan Asumsi Penelitian}
Hipotesis dari tulisan ini adalah
\begin{enumerate}
    \item Data penelitian ini diambil dari data NASA yang dianggap sudah melalui proses validasi dan tidak mengandung kesalahan pencatatan yang signifikan ada di wilayah Surabaya
    \item Data curah hujan yang digunakan adalah data harian selama 2015 sampai 2025
    \item Model Spasio temporal yang digunakan adalah STARMA yang dianggap mampu merepresentasikan spatial dan temporal dalam curah hujan tanpa penggabungan model hybrid atau \textit{machine learning}
    \item Data NASA berupa tangkapan satelit di wilayah Surabaya didasarkan pada koordinat geografis (longitude dan latitude).
\end{enumerate}